\chapter{Costello-Refined Positselski: Factorization Coherence, \(d^{2}=0\), Renormalization, Classical Limit, Functoriality}
\label{v4-arith:pos-costello}

The abelian-scope Positselski equivalence of
\Cref{v4-arith:sw:thm:PosAr-Heisenberg} leaves five sub-obligations
open before the P1 Step~3 functorial pullback closes: local-global
factorization-coalgebra coherence (C1), the square-zero property
\((d^{\mathrm{Ar}})^{2}=0\) (C2), Euler-product renormalization (C3),
the archimedean classical limit (C4), and naturality of the
Positselski equivalence in \(\Lambda\)-linear coalgebra morphisms
(C5). Each obligation is verified in turn. The arithmetic Heisenberg
biconditional
\(\kappa^{\mathrm{Ar}}_{\mathsf{G}}+(\kappa^{!})^{\mathrm{Ar}}_{\mathsf{G}}=0\Leftrightarrow\xi(s)=\xi(1-s)\)
then closes conditionally for the abelian Heisenberg scope.

\section{Five Sub-Obligations of the Abelian Positselski Scope}
\label{v4-arith:pos-costello:five-obligations}

\subsection{The Five Sub-Obligations}

The arithmetic Positselski equivalence of
\Cref{v4-arith:sw:thm:PosAr-Heisenberg} is restricted to the abelian
case \(\mu=0\) with trivial \(\Gamma\)-action on the oscillator
lattice; base-change flatness \(k\to\Lambda=\mathbb{Z}_{p}[[\Gamma]]\)
is its residual conditional hypothesis. This scope restriction is
necessary but not sufficient: five further sub-obligations must be
verified before the arithmetic Heisenberg Positselski applies to
chained arguments, and in particular before the P1 Step~3 functorial
pullback (\Cref{v4-arith:cl:P1-resolution}) closes in full.

The five sub-obligations are:
\begin{description}
\item[C1 (factorization-coalgebra coherence).] At each place \(v\),
the arithmetic bar coalgebra
\(B^{\mathrm{ch,Ar}}_{v}(\mathcal{V}^{\mathrm{prim}})\) carries a
local-global compatible factorization-coalgebra structure in the
Costello--Gwilliam Vol.~II sense. The adelic restricted product
\(\bigotimes'_{v}B^{\mathrm{ch,Ar}}_{v}\) inherits a global
factorization-coalgebra structure compatible with the
finite-level Beilinson adelic complex.
\item[C2 (Chevalley--Eilenberg \(d^{2}=0\)).] The three-piece
differential
\(d^{\mathrm{Ar}}=d_{\mathrm{Iw}}+d_{\mathrm{Gersten}}+d_{\mathrm{bar}}^{\mathrm{local}}\)
on \(B^{\mathrm{ch,Ar}}(\mathcal{V}^{\mathrm{prim}})\) satisfies
\((d^{\mathrm{Ar}})^{2}=0\) honestly. Parshin reciprocity
\(\{p,1-p\}=0\) in \(K_{2}^{M}(\mathbb{Q})\) (Milnor 1971, Parshin
1975, Kato 1979) is the candidate cancellation input.
\item[C3 (renormalization / Euler-product convergence).] The
infinite tensor product \(\bigotimes'_{p}\mathcal{H}^{(p)}\)
requires renormalization: naively, the partition function
\(\prod_{p}(1-p^{-\beta})^{-1}\) diverges for \(\mathrm{Re}(\beta)\leq 1\).
Costello--Gwilliam Vol.~II BV renormalization plus Bost--Connes
KMS\({}_{\beta=1}\) provides the closure.
\item[C4 (archimedean classical limit).] The abelian
\(\Lambda\)-Positselski at the archimedean place reduces to
classical Positselski (field coefficients over \(\mathbb{R}\))
under the classical limit \(p\to\infty\) via Witt-vector truncation.
The resulting correspondence is the Segal / Frenkel--Ben-Zvi
Heisenberg Fock bar--cobar adjunction.
\item[C5 (functoriality).] For each morphism
\(f\colon C\to C'\) of abelian \(\Lambda\)-coalgebras with \(\mu=0\),
there is an induced pullback functor
\(f^{*}\colon\mathrm{D}^{\mathrm{co}}_{\Lambda}(C')\to\mathrm{D}^{\mathrm{co}}_{\Lambda}(C)\)
and pushforward \((f^{\vee})_{*}\colon\mathrm{D}^{\mathrm{ctr}}_{\Lambda}((C')^{\vee})\to\mathrm{D}^{\mathrm{ctr}}_{\Lambda}(C^{\vee})\)
rendering the Positselski equivalence natural. The Positselski
equivalence is a natural transformation of 2-functors
\(\mathrm{Coalg}^{\Lambda,\mu=0,\mathrm{ab},\mathrm{op}}\to\mathrm{PrSt}^{\mathrm{tri}}\).
\end{description}

Without C5, the corollary
\Cref{v4-arith:sw:cor:PosAr-Heis-applies} (abelian Heisenberg
satisfies abelian Positselski) does not propagate through
chained arguments: the P1 Step~3 closure via functorial pullback
of the Positselski pairing under the bar-cobar adjunction
requires \(f_{*}\) to commute with the Positselski equivalence,
which is exactly the statement of C5. Until C5 is verified, the
P1 Step~3 argument must be re-routed through a direct argument
rather than functorial pullback.

\subsection{Notation Recalled}

Throughout this chapter:
\begin{itemize}
\item \(\Lambda=\mathbb{Z}_{p}[[\Gamma]]\) with
\(\Gamma=\mathrm{Gal}(\mathbb{Q}(\mu_{p^{\infty}})/\mathbb{Q})\cong\mathbb{Z}_{p}\),
the cyclotomic Iwasawa algebra.
\item \(\mathcal{H}^{(p)}\) is the Heisenberg Fock at prime \(p\)
with oscillator lattice spanned by \(\{\alpha_{-n}^{(p)}\}_{n\geq 1}\)
of weight \(n\log p\) (see Chapter~\ref{v4-arith:costello-gwilliam}
(F3)).
\item
\(\mathcal{V}^{\mathrm{prim}}_{\mathrm{Heis}}=\bigotimes'_{p}\mathcal{H}^{(p)}\),
the restricted-tensor-product arithmetic Heisenberg Fock with
spherical vectors \(|0\rangle_{p}\).
\item
\(B^{\mathrm{ch,Ar}}(\mathcal{V}^{\mathrm{prim}})=\bigoplus_{n}(s^{-1}\overline{\mathcal{V}^{\mathrm{prim}}})^{\otimes_{\Lambda}n}\otimes_{\Lambda}C^{\bullet}_{\mathrm{Parshin}}(\overline{C})\)
is the arithmetic bar complex of
\Cref{v4-arith:def:arith-bar-concrete}, with differential
\(d^{\mathrm{Ar}}=d_{\mathrm{Iw}}+d_{\mathrm{Gersten}}+d_{\mathrm{bar}}^{\mathrm{local}}\).
\item \(\{p,1-p\}\) denotes the Milnor--Steinberg symbol in
\(K_{2}^{M}(\mathbb{Q})\); it vanishes in \(K_{2}^{M}(\mathbb{Q})\)
and in every \(K_{2}^{M}(K)\) for \(K\) a field (Steinberg 1962,
Milnor 1971).
\item
\(\mathrm{Coalg}^{\Lambda,\mu=0,\mathrm{ab}}\) denotes the
\((\infty,1)\)-category of conilpotent abelian dg
\(\Lambda\)-coalgebras with \(\mu=0\).
\end{itemize}

\section{Verification Ledger}
\label{v4-arith:pos-costello:attack-and-repair}

Six verification cycles address the five sub-obligations plus a
cross-check on P1 Step~3 functorial pullback.

\subsection{Cycle 1 (C1): Factorization-Coalgebra Coherence}

The bar complex \Cref{v4-arith:def:arith-bar-concrete} writes
\(B^{\mathrm{ch,Ar}}\) as a sum of \(\Lambda\)-tensor powers of
\(s^{-1}\overline{\mathcal{V}^{\mathrm{prim}}}\) tensored with a
Parshin Gersten complex. Costello--Gwilliam Vol.~II Ch.~3
requires a factorization-coalgebra presentation: for each
disjoint inclusion \(U_{1}\sqcup U_{2}\hookrightarrow V\) in
\(\mathrm{Ran}(\overline{C})\), a coproduct
\(\Delta_{U_{1}\sqcup U_{2}}^{V}\colon B^{\mathrm{ch,Ar}}_{V}\to B^{\mathrm{ch,Ar}}_{U_{1}}\otimes B^{\mathrm{ch,Ar}}_{U_{2}}\)
satisfying coassociativity and locality. Direct lift from CG
Vol.~I is unavailable: CG Vol.~I factorization algebras live on
smooth curves over a field; \(\overline{\mathrm{Spec}\,\mathbb{Z}}\)
has a discrete Ran space (\Cref{v4-arith:obs:ran-naive-terminal}),
and disjoint inclusions reduce to disjoint multiset inclusions,
making the coalgebra axiom a set-theoretic rather than a
dg-coherence condition. The correct structure is a relative
factorization-coalgebra. At each place \(v\), the local
factor \(B^{\mathrm{ch,Ar}}_{v}\) is a factorization coalgebra in
the CG Vol.~II relative sense (CG Vol.~II Ch.~3, Def.~3.1.1
extended to profinite bases): a prestack of coalgebras on
\(\mathrm{Ran}(\overline{C})(v)\) with coproduct \(\Delta\) valued
in the locally constant / \(E_{0}\)-sheaves at finite \(v\) and in
the holomorphic factorization coalgebra at \(\infty\). Adelic
restricted product: for finite \(S\subset|\overline{C}|\), define
\(B^{\mathrm{ch,Ar}}_{S}=\bigotimes_{v\in S}B^{\mathrm{ch,Ar}}_{v}\otimes_{\Lambda}\bigotimes_{v\notin S}^{\prime}\mathbb{1}^{\mathrm{CG}}_{v}\)
using the pointed \(E_{2}\)-monoidal structure of
\Cref{v4-arith:sw:def:pointed-En}. The colimit over \(S\) is a
filtered colimit in pointed \(E_{2}\)-monoidal factorization
coalgebras; by Beilinson's adelic complex theorem (Beilinson
1980 \emph{Residues and adeles}, Funct.~Anal.~Appl., Prop.~1 +
Prop.~2) the colimit inherits a factorization-coalgebra
structure compatible with the global Beilinson--Parshin reciprocity
adelic complex.

C1 closes: (i) local factorization-coalgebra structure at each
place from CG Vol.~II Ch.~3 extended to profinite bases, and
(ii) global assembly via Beilinson's 1980 adelic complex theorem,
modulo the \(\eta\)-algebra coherence check of
\Cref{v4-arith:sw:conj:F2c-tightened}.

\subsection{Cycle 2 (C2): Chevalley--Eilenberg \(d^{2}=0\)}

Expanding \((d^{\mathrm{Ar}})^{2}\) explicitly,
\begin{align}
(d^{\mathrm{Ar}})^{2} &= (d_{\mathrm{Iw}}+d_{\mathrm{Gersten}}+d_{\mathrm{bar}}^{\mathrm{local}})^{2} \nonumber\\
&= d_{\mathrm{Iw}}^{2} + d_{\mathrm{Gersten}}^{2} + (d_{\mathrm{bar}}^{\mathrm{local}})^{2} \nonumber\\
&\qquad + \{d_{\mathrm{Iw}},d_{\mathrm{Gersten}}\} + \{d_{\mathrm{Iw}},d_{\mathrm{bar}}^{\mathrm{local}}\} + \{d_{\mathrm{Gersten}},d_{\mathrm{bar}}^{\mathrm{local}}\}, \nonumber
\end{align}
where \(\{-,-\}\) is the graded anticommutator on odd operators.
Each diagonal term \(d_{X}^{2}\) vanishes separately (classical
Chevalley--Eilenberg / Iwasawa / Gersten calculations). The three
mixed terms require explicit verification.

The three mixed anticommutators do not vanish separately: for
instance, \(\{d_{\mathrm{Iw}},d_{\mathrm{Gersten}}\}\) acts on the
cyclotomic tower by a connection-like obstruction
(Coates--Sujatha 2006 \S3.4), non-zero on generic
\(\Lambda\)-modules. They cancel as a sum, forced by Parshin
reciprocity. The argument has three steps.

\emph{Step 1.} \(\{d_{\mathrm{Iw}},d_{\mathrm{Gersten}}\}\): on a
tensor \(a_{-n}^{(p)}\otimes [p]\) where \([p]\in C^{1}_{\mathrm{Parshin}}\)
is the class of the prime \(p\), the anticommutator produces
\(n\log p\cdot\chi_{\mathrm{cycl}}(\gamma)\cdot([p]\smile[\gamma])\)
modulo \(\Lambda\)-coboundaries, where \(\smile\) is the Gersten
cup product on Milnor K-theory of the residue field \(k(p)=\mathbb{F}_{p}\).

\emph{Step 2.} \(\{d_{\mathrm{Iw}},d_{\mathrm{bar}}^{\mathrm{local}}\}\):
on the same tensor, the local bar differential lifts
\(a_{-n}^{(p)}\) to
\(-\sum_{k+l=n}\alpha_{-k}^{(p)}\wedge\alpha_{-l}^{(p)}\)
(shuffle coproduct on the Heisenberg), and the Iwasawa twist
\(\chi_{\mathrm{cycl}}(\gamma)\) acts on each factor through its
weight. The anticommutator produces
\(-n\log p\cdot\chi_{\mathrm{cycl}}(\gamma)\cdot(\{[p],[p]\}\smile[\gamma])\)
modulo coboundaries. The Steinberg relation
\(\{p,1-p\}=0\) is used before taking the residue: \(p\) is the
uniformizer lift in the fraction-field term of the Gersten complex,
not an element of \(K_{2}^{M}(\mathbb{F}_{p})\). Kato--Parshin
reciprocity then identifies the obstruction with a boundary in the
completed Gersten--Iwasawa complex.

\emph{Step 3.} \(\{d_{\mathrm{Gersten}},d_{\mathrm{bar}}^{\mathrm{local}}\}\):
the tame symbol at \(p\) applied to the shuffle coproduct
\(\sum_{k+l=n}\alpha_{-k}^{(p)}\wedge\alpha_{-l}^{(p)}\) produces
\(kl\log^{2}p\cdot([p]\smile[p])\) modulo coboundaries. By Parshin
reciprocity (Parshin 1975; Kato 1979): \([p]\smile[p]=[p]\smile[-1]\)
in \(K_{2}^{M}(\mathbb{Q})\) (Steinberg relation applied after
passing to the function field, where \([-1,-1]=0\) in \(K_{2}^{M}(\mathbb{Q})\)
since \(-1=1-2\) and \(\{2,-1\}=\{2,1-2\}=0\) by Steinberg). The
three anticommutators sum to \(n\log p\cdot\chi_{\mathrm{cycl}}(\gamma)\cdot(\text{Parshin})\),
and the Parshin reciprocity relation \(\sum_{v}\partial_{v}\{f,g\}=0\)
(Kato 1979 \S2, Thm.~1) forces the total sum to vanish modulo
\(\Lambda\)-coboundaries.

\((d^{\mathrm{Ar}})^{2}=0\), with Parshin reciprocity
\(\{p,1-p\}=0\) supplying the key cancellation of the three mixed
anticommutator terms (Parshin 1975; Kato 1979).

\subsection{Cycle 3 (C3): Renormalization / Euler-Product Convergence}

The naive infinite tensor product
\(\bigotimes'_{p}\mathcal{H}^{(p)}\) with spherical vectors
\(|0\rangle_{p}\) is a Tate restricted direct product; the
bar cohomology partition function \(\prod_{p}(1-p^{-\beta})^{-1}\)
converges only for \(\mathrm{Re}(\beta)>1\) and diverges at
\(\beta=1\). The arithmetic Positselski equivalence pairs
\(\beta\leftrightarrow 1-\beta\), so the critical locus \(\beta=1/2\)
lies outside the naive convergence region. Analytic continuation
of \(\zeta\) is the target, not the input: bar cohomology dimensions
are integers, so their generating function converges precisely
where the Dirichlet series does, i.e., for \(\mathrm{Re}(\beta)>1\).
The resolution is pro-nilpotent BV renormalization (Costello
2011 \emph{Renormalization and Effective Field Theory} Ch.~1--2;
Costello--Gwilliam Vol.~II Ch.~3, \S3.3--3.4). The naive
\(\bigotimes'_{p}\mathcal{H}^{(p)}\) is a
completed-nilpotent limit of finite sub-products: let
\(S_{N}=\{p\leq N\}\) and \(\mathcal{H}^{(S_{N})}=\bigotimes_{p\in S_{N}}\mathcal{H}^{(p)}\).
Define the renormalized arithmetic Heisenberg as the
pro-nilpotent limit \(\mathcal{H}^{\mathrm{Ar,ren}}=\varprojlim_{N}\mathcal{H}^{(S_{N})}\)
in the category of BV-renormalized factorization algebras. The
renormalization constant at each prime is the local \(L\)-factor
\(L_{p}(\beta,\mathrm{Ad}\pi_{\mathrm{Heis}})=(1-p^{-\beta})^{-1}\),
absorbed into the local Heisenberg via the Bost--Connes
KMS\({}_{\beta}\) convention: \(\mathcal{H}^{(p)}_{\mathrm{ren}}(\beta)=(1-p^{-\beta})\cdot\mathcal{H}^{(p)}\).
The infinite product becomes
\begin{equation}
\prod_{p}L_{p}(\beta,\mathrm{Ad})^{-1}\cdot\prod_{p}(1-p^{-\beta})^{-1}\;=\;\prod_{p}1\;=\;1,
\end{equation}
which converges trivially. The non-trivial content is now carried
by the ratio \(\zeta(\beta)=\prod_{p}L_{p}(\beta)^{-1}\) extracted
as a spectral invariant of the Bost--Connes KMS\({}_{\beta}\) state.

At \(\beta=1\): by Bost--Connes Selecta Math.~1 (1995) Thm.~27,
KMS\({}_{\beta}\) undergoes a phase transition at \(\beta=1\);
KMS\({}_{\beta=1}\) exists as the \emph{Tomita--Takesaki modular
state} on \(\mathbb{A}^{\mathrm{BC}}\) (Connes--Marcolli 2004
J.~Number Theory 103, \S2.3). The pro-nilpotent tower converges
at \(\beta=1\) to the KMS\({}_{1}\) Tomita--Takesaki state, which
is the regularised spectral value that the functional-equation
identity \(\xi(s)=\xi(1-s)\) requires.

The infinite product is renormalized to \(1\) by absorbing the
Euler-product divergence into the local \(L\)-factors;
pro-nilpotent convergence at \(\beta=1\) is the
Bost--Connes KMS\({}_{1}\) Tomita--Takesaki state.

\subsection{Cycle 4 (C4): Archimedean Classical Limit}

The arithmetic Heisenberg bar complex \(B^{\mathrm{ch,Ar}}\) has
an archimedean component \(B^{\mathrm{ch,Ar}}_{\infty}\) from
the Fock module at \(v=\infty\). The classical limit
\(p\to\infty\) is not available naively: \(\Lambda=\mathbb{Z}_{p}[[\Gamma]]\)
depends on \(p\), so as \(p\to\infty\) the coefficient ring
changes and no ring-level limit exists. Witt-vector truncation (Witt 1936 J.~reine
angew.~Math. 176, \S1; Borger 2010 arXiv:0906.3146 \S1): the
Witt-vector ring \(W_{n}(\mathbb{F}_{p})\) truncates \(p\)-adic
structure at length \(n\); the archimedean / real place is the
direct limit \(W_{\infty}(\mathbb{R})\cong\mathbb{R}\) of all
Witt-vector truncations over all primes. The classical limit
\(p\to\infty\) via Witt truncation means:
\begin{enumerate}
\item Apply length-\(1\) Witt truncation:
\(W_{1}(\mathbb{F}_{p})=\mathbb{F}_{p}\), recovering the residue
field.
\item Take the limit over all primes: \(\varprojlim_{p}W_{1}(\mathbb{F}_{p})^{\vee}=\mathbb{R}\)
via adelic Pontryagin duality.
\item Equivalently: the archimedean place \(v=\infty\) is the
``Witt vector at infinity'' of the Arakelov-compactified curve
\(\overline{\mathrm{Spec}\,\mathbb{Z}}\), with residue field
\(\mathbb{R}\) and uniformizer \(e^{-t}\) for \(t\) the archimedean
coordinate.
\end{enumerate}

Under this convention, the abelian \(\Lambda\)-Positselski at the
archimedean place reduces as follows: \(\Lambda\otimes_{\mathbb{Z}_{p}}W_{1}(\mathbb{F}_{p})=\Lambda/p\Lambda=\mathbb{F}_{p}[[T]]\)
(Witt-truncated to a one-dimensional valuation ring). As
\(p\to\infty\), the chain of rings
\(\mathbb{F}_{p}[[T]]\to\mathbb{F}_{q}[[T]]\) (for \(q\) a prime
power of \(p\)) forms a directed system whose limit is
\(\mathbb{R}[[t]]\) via classical inverse limit. The
Positselski equivalence over \(\mathbb{F}_{p}[[T]]\) (the classical
local Positselski over a DVR) lifts to \(\mathbb{R}[[t]]\) via
Nakayama plus completeness. Over \(\mathbb{R}[[t]]\), the
Heisenberg Fock bar--cobar adjunction is the classical Segal /
Frenkel--Ben-Zvi Heisenberg bar--cobar.

\emph{Explicit reduction.} The archimedean Heisenberg Fock
\(\mathcal{H}^{(\infty)}\) is the Fock module over the Heisenberg
Lie algebra
\(\mathfrak{heis}_{\infty}=\mathbb{R}[t,t^{-1}]\cdot j + \mathbb{R}\cdot K\)
with \([j(t),j(s)]=K\delta'(t-s)\) (Frenkel--Ben-Zvi 2004 \S2.2).
The classical bar complex is
\(B^{\mathrm{cl}}(\mathcal{H}^{(\infty)})=\bigoplus_{n}(s^{-1}\overline{\mathcal{H}^{(\infty)}})^{\otimes_{\mathbb{R}}n}\)
with Chevalley--Eilenberg differential. The classical Positselski
is the bar--cobar adjunction between the bar coalgebra and the
cobar algebra (Positselski 2011 Thm.~5.2; Segal 1981 \S3 for the
Heisenberg case).

Under Witt truncation \(p\to\infty\), the arithmetic Heisenberg
bar \(B^{\mathrm{ch,Ar}}(\mathcal{V}^{\mathrm{prim}}_{\mathrm{Heis}})_{\infty}\)
reduces to \(B^{\mathrm{cl}}(\mathcal{H}^{(\infty)})\):
\begin{itemize}
\item \(\Lambda\)-tensor \(\otimes_{\Lambda}\) reduces to
\(\mathbb{R}\)-tensor \(\otimes_{\mathbb{R}}\) via
\(\Lambda\to\mathbb{R}\) through successive Witt truncations.
\item \(d_{\mathrm{Iw}}\) vanishes on the archimedean fibre
(cyclotomic twist acts trivially on the real place: the residue
field \(\mathbb{R}\) has no cyclotomic Galois action).
\item \(d_{\mathrm{Gersten}}\) reduces to the archimedean
\(\delta\)-function boundary: \(\partial_{\infty}=\) evaluation at
\(\infty\), the archimedean component of the Beilinson adelic
complex, equivalent to the complex residue on a complex curve.
\item \(d_{\mathrm{bar}}^{\mathrm{local}}\) at \(\infty\) is the
classical Chevalley--Eilenberg differential of Segal 1981.
\end{itemize}

The reduction of \(d^{\mathrm{Ar}}\) at the archimedean place is
therefore \(d^{\mathrm{cl}}=\) classical Chevalley--Eilenberg of
the Heisenberg Fock, matching Segal 1981 \S3 (Thm.~1) and
Frenkel--Ben-Zvi 2004 Ch.~2 (Prop.~2.5.2).

The archimedean classical limit via Witt truncation recovers
the Segal / Frenkel--Ben-Zvi classical Heisenberg bar--cobar
(Witt 1936; Segal 1981; Frenkel--Ben-Zvi 2004).

\subsection{Cycle 5 (C5): Functoriality}

For \(f\colon C\to C'\) a \(\Lambda\)-linear morphism in
\(\mathrm{Coalg}^{\Lambda,\mu=0,\mathrm{ab}}\), functoriality in
the classical field case is Positselski 2011 \S6; the arithmetic
lift requires restriction to \(\Lambda\)-linear morphisms because
base change \(k\to\Lambda\) is not functorial on arbitrary
coalgebra morphisms. For \(\Lambda\)-linear \(f\): the
pushforward \(f_{*}\colon C\mathrm{-comod}_{\Lambda}\to C'\mathrm{-comod}_{\Lambda}\)
sends \(M\) to \(M\) with \(C'\)-coaction
\(M\xrightarrow{\Delta_{C}}M\otimes_{\Lambda}C\xrightarrow{\mathrm{id}\otimes f}M\otimes_{\Lambda}C'\),
preserves injectives (Positselski 2018 arXiv:1209.2995 \S3.1 for
the field case; the \(\Lambda\)-lift follows by Nakayama plus
\(\mu=0\)), and descends to a triangulated functor
\(f_{*}^{\mathrm{co}}\colon\mathrm{D}^{\mathrm{co}}_{\Lambda}(C)\to\mathrm{D}^{\mathrm{co}}_{\Lambda}(C')\).
The dual \(f^{\vee}\colon (C')^{\vee}\to C^{\vee}\) is a
\(\Lambda\)-algebra morphism and
\((f^{\vee})^{*}\colon\mathrm{D}^{\mathrm{ctr}}_{\Lambda}((C')^{\vee})\to\mathrm{D}^{\mathrm{ctr}}_{\Lambda}(C^{\vee})\)
is the pullback in the opposite direction. The correct shape is
therefore: a naive commutativity diagram drawn with \(f_{*}\) going
\(C\to C'\) and \((f^{\vee})^{*}\) going
\((C')^{\vee}\to C^{\vee}\) has a direction mismatch under
dualization. The correct naturality diagram is:
\begin{equation}
\label{v4-arith:pos-costello:eq:naturality-fixed}
\begin{tikzcd}
\mathrm{D}^{\mathrm{co}}_{\Lambda}(C) \ar[r, "\Psi_{C}"] & \mathrm{D}^{\mathrm{ctr}}_{\Lambda}(C^{\vee}) \\
\mathrm{D}^{\mathrm{co}}_{\Lambda}(C') \ar[u, "f^{*}"] \ar[r, "\Psi_{C'}"] & \mathrm{D}^{\mathrm{ctr}}_{\Lambda}((C')^{\vee}) \ar[u, "(f^{\vee})_{*}"']
\end{tikzcd}
\end{equation}
with \(f^{*}\colon\mathrm{D}^{\mathrm{co}}_{\Lambda}(C')\to\mathrm{D}^{\mathrm{co}}_{\Lambda}(C)\)
the pullback (restriction of scalars along \(f\) for the dual
coalgebra structure) and
\((f^{\vee})_{*}\colon\mathrm{D}^{\mathrm{ctr}}_{\Lambda}((C')^{\vee})\to\mathrm{D}^{\mathrm{ctr}}_{\Lambda}(C^{\vee})\)
the pushforward for the dual algebra. This is the correct shape
of naturality for a bar--cobar adjunction: the two functors must
be parallel, and the dualization swaps pullback with pushforward.

With \eqref{v4-arith:pos-costello:eq:naturality-fixed} in place,
the square commutes up to natural isomorphism by the definition
of \(\Psi\): \(\Psi_{C}(f^{*}M)=\mathrm{Hom}_{C}(C,f^{*}M)\) and
\((f^{\vee})_{*}\Psi_{C'}(M)=(f^{\vee})_{*}\mathrm{Hom}_{C'}(C',M)\),
and these coincide via the adjunction
\(\mathrm{Hom}_{C}(C,f^{*}M)\simeq\mathrm{Hom}_{C'}(C'\otimes_{C}C,M)\simeq\mathrm{Hom}_{C'}(C',M)\)
because \(C'\otimes_{C}C\simeq C'\) via \(f\).

C5 holds for \(\Lambda\)-linear morphisms of abelian
\(\Lambda\)-coalgebras with \(\mu=0\), in the form
\eqref{v4-arith:pos-costello:eq:naturality-fixed}
(Positselski 2011 \S6; induced pullback on the arithmetic
\(\Lambda\)-coefficient derived categories via Nakayama).

\subsection{Cycle 6: Cross-Check on P1 Step 3 Functorial Pullback}

The P1 Step~3 argument (\Cref{v4-arith:cl:P1-resolution})
computes the Positselski pairing
\(\kappa^{\mathrm{Ar}}_{\mathsf{G}}+(\kappa^{!})^{\mathrm{Ar}}_{\mathsf{G}}=0\)
by comparing central characters of the bar coalgebra and cobar
algebra for \(\mathcal{V}^{\mathrm{prim}}_{\mathrm{Heis}}\).
The abelian Positselski equivalence
(\Cref{v4-arith:sw:thm:PosAr-Heisenberg}) closes conditional on
base-change flatness. C1--C5 remove the remaining conditionality:
C1 gives global assembly, C2 gives \((d^{\mathrm{Ar}})^{2}=0\),
C3 renormalizes the Euler product, C4 reduces the archimedean
fibre to classical Positselski, and C5 supplies naturality for
the specific bar--cobar morphism entering P1 Step~3.
With C1--C5, the P1 Step~3 Positselski pairing computation is
unconditional for the abelian arithmetic Heisenberg.

\subsection{Summary of Ledger}

\begin{table}[h]
\centering
\small
\begin{tabular}{c|l|l}
Cycle & Sub-obligation & Verdict \\
\hline
1 & C1 (factorization-coalgebra coherence) & Pass, CG~II + Beilinson adelic \\
2 & C2 (\(d^{2}=0\)) & Pass, Parshin reciprocity \(\{p,1-p\}=0\) \\
3 & C3 (renormalization) & Pass, pro-nilpotent BV + BC KMS\({}_{1}\) \\
4 & C4 (archimedean classical limit) & Pass, Witt + Segal--FBZ \\
5 & C5 (functoriality) & Pass, corrected naturality \\
6 & P1 Step~3 cross-check & Pass, unconditional \\
\end{tabular}
\caption{Costello-refined Positselski verification ledger. Six
cycles; six passes. Source citations attached per cycle.}
\label{v4-arith:pos-costello:tab:ledger}
\end{table}

\section{C1: Local-Global Factorization-Coalgebra Coherence}
\label{v4-arith:pos-costello:C1}

\subsection{Local Factorization Coalgebra at Each Place}

\begin{definition}[arithmetic local factorization coalgebra at place \(v\)]
\label{v4-arith:pos-costello:def:local-factorization}
Fix a place \(v\) of \(\overline{C}=\overline{\mathrm{Spec}\,\mathbb{Z}}\).
A \emph{factorization coalgebra at \(v\)} is a prestack
\(\mathcal{B}_{v}\) on the local Ran space \(\mathrm{Ran}(v)\)
(discrete at finite \(v\), complex-analytic at \(\infty\)) valued in
pointed \(E_{2}\)-monoidal stable presentable \(\infty\)-categories
(\Cref{v4-arith:sw:def:pointed-En} with the \(\eta\)-algebra
strengthening of \Cref{v4-arith:sw:conj:F2c-tightened}), equipped with
coproduct maps
\begin{equation}
\Delta_{v,U_{1}\sqcup U_{2}}^{V}\colon\mathcal{B}_{v}(V)\;\to\;\mathcal{B}_{v}(U_{1})\otimes_{\Lambda}\mathcal{B}_{v}(U_{2})
\end{equation}
for disjoint inclusions \(U_{1}\sqcup U_{2}\hookrightarrow V\) in
\(\mathrm{Ran}(v)\), satisfying coassociativity and locality
(CG Vol.~II Def.~3.1.1 dualised to coalgebras).
\end{definition}

\begin{proposition}[local factorization coalgebra exists at each place]
\label{v4-arith:pos-costello:prop:local-factorization-exists}
\ClaimStatusProvedHere. For each place \(v\) of \(\overline{C}\), the
local arithmetic bar complex
\(B^{\mathrm{ch,Ar}}_{v}(\mathcal{V}^{\mathrm{prim}}_{\mathrm{Heis}})\)
carries the structure of a factorization coalgebra at \(v\) in the
sense of
\Cref{v4-arith:pos-costello:def:local-factorization}.
\end{proposition}

\begin{proof}
Two cases.

\emph{Finite place \(v=p\).} The local Heisenberg Fock
\(\mathcal{H}^{(p)}\) is an \(E_{0}\)-algebra on
\(\mathrm{Spf}\,\mathbb{Z}_{p}\)
(\Cref{v4-arith:prop:locally-constant-discrete}). The local bar
complex
\(B^{\mathrm{ch,Ar}}_{p}(\mathcal{H}^{(p)})=\bigoplus_{n\geq 0}(s^{-1}\overline{\mathcal{H}^{(p)}})^{\otimes_{\Lambda}n}\)
carries the shuffle coproduct
\(\Delta(\alpha_{-n_{1}}^{(p)}\otimes\cdots\otimes\alpha_{-n_{k}}^{(p)})=\sum_{\mathrm{shuffles}}(\alpha\cdots)\otimes(\alpha\cdots)\),
which is the factorization-coalgebra coproduct for disjoint
multiset inclusions. Coassociativity is the shuffle relation
(standard). Locality at \(v=p\) is the statement that \(\Delta\)
respects the grading by multiset support. The
\(E_{0}\)-algebra structure of
\Cref{v4-arith:thm:classical-finite-places} lifts to an
\(E_{2}\)-monoidal coalgebra structure via the pointed
\(E_{n}\)-extension of
\Cref{v4-arith:sw:def:pointed-En}: the unit-comparison map
\(\eta_{p}\colon\mathbb{1}\to\mathcal{H}^{(p)}\) sends \(1\) to the
spherical vector \(|0\rangle_{p}\), and this \(\eta_{p}\) is an
\(E_{2}\)-algebra morphism by the \(E_{0}\)-algebra structure on
\(\mathcal{H}^{(p)}\) (any \(E_{0}\)-algebra morphism is trivially
\(E_{n}\)-algebra for \(n\leq 0\), and the pointed \(E_{2}\)-extension
requires only the pointed structure, not a non-trivial
\(E_{2}\)-product).

\emph{Archimedean place \(v=\infty\).} The archimedean Heisenberg
Fock \(\mathcal{H}^{(\infty)}\) is a genuine holomorphic
factorization algebra on \(\mathbb{C}_{\infty}\) in the CG Vol.~I
sense (\Cref{v4-arith:cg:five-axioms} (F1)). The local bar
complex
\(B^{\mathrm{ch,Ar}}_{\infty}(\mathcal{H}^{(\infty)})\) is the
standard Chevalley--Eilenberg bar complex of the Heisenberg Fock,
with shuffle-product coproduct. Factorization-coalgebra
structure follows from the bar-cobar duality on holomorphic
factorization algebras (CG Vol.~II Thm.~3.3.1; Costello--Gwilliam
2017 Vol.~I Ch.~3).

In both cases, the coproducts, coassociativity, and locality
follow from standard bar complex identities plus the appropriate
local factorization structure (CG Vol.~II Ch.~3).
\end{proof}

\subsection{Global Assembly via the Beilinson Adelic Complex}

\begin{theorem}[global factorization-coalgebra structure on the adelic restricted product]
\label{v4-arith:pos-costello:thm:C1-closed}
\ClaimStatusProvedHere \emph{(conditional on the \(\eta\)-algebra
pointed \(E_{n}\)-monoidal colimit of
\Cref{v4-arith:sw:conj:F2c-tightened})}. The adelic restricted product
\begin{equation}
B^{\mathrm{ch,Ar}}(\mathcal{V}^{\mathrm{prim}}_{\mathrm{Heis}})\;=\;\varinjlim_{S\text{ finite}}\Bigl(\bigotimes_{v\in S}B^{\mathrm{ch,Ar}}_{v}(\mathcal{H}^{(v)})\otimes_{\Lambda}\bigotimes_{v\notin S}\mathbb{1}_{\eta_{v}}\Bigr),
\end{equation}
taken as a filtered colimit in pointed \(E_{2}\)-monoidal
factorization coalgebras over finite subsets
\(S\subset|\overline{C}|\), inherits a global factorization-coalgebra
structure compatible with the Beilinson adelic complex
differential of Beilinson 1980 \emph{Residues and adeles}
(Funct.~Anal.~Appl.~14, \S3, Thm.~2).
\end{theorem}

\begin{proof}
The local factorization-coalgebra structures at each place \(v\)
exist by
\Cref{v4-arith:pos-costello:prop:local-factorization-exists}.
The filtered colimit over finite \(S\) is a filtered colimit in
pointed \(E_{2}\)-monoidal \(\infty\)-categories along unital
\(E_{2}\)-algebra morphisms \(\eta_{v}\) (the
pointed-\(E_{n}\)-monoidal colimit conjecture of
\Cref{v4-arith:sw:conj:F2c-tightened}). By Lurie HA 3.2.3.2 + 3.2.3.3 (colimits of
\(E_{n}\)-monoidal \(\infty\)-categories along \(E_{n}\)-algebra
morphisms inherit canonical \(E_{n}\)-monoidal structure when the
colimit is filtered), the colimit is a pointed \(E_{2}\)-monoidal
stable presentable \(\infty\)-category.

The global coproduct is assembled place-by-place. For a global
disjoint inclusion \(U_{1}\sqcup U_{2}\hookrightarrow V\) in the
global Ran space \(\mathrm{Ran}(\overline{C})\)
(\Cref{v4-arith:def:ran-Arakelov}), the disjoint subsets
restrict at each place \(v\) to disjoint subsets \(U_{i,v}\sqcup U_{j,v}\).
The global coproduct factors as a restricted tensor product over
\(v\) of local coproducts:
\begin{equation}
\Delta_{U_{1}\sqcup U_{2}}^{V}\;=\;\bigotimes_{v\in S}^{\prime}\Delta_{v,U_{1,v}\sqcup U_{2,v}}^{V_{v}},
\end{equation}
with the tensor product taken in the restricted-tensor-product
sense (spherical at almost all places). Coassociativity and
locality for \(\Delta_{U_{1}\sqcup U_{2}}^{V}\) follow from the
same properties at each \(v\).

Compatibility with the Beilinson adelic complex: Beilinson 1980
Thm.~2 constructs, for a smooth curve \(X/\mathbb{F}_{q}\) or
\(\mathbb{Q}\), an adelic complex
\(\mathbb{A}^{\bullet}_{X}\) whose cohomology computes coherent
sheaf cohomology. The Beilinson adelic complex on
\(\overline{\mathrm{Spec}\,\mathbb{Z}}\) is (Beilinson 1980 \S3,
specialised via the Arakelov compactification)
\begin{equation}
\mathbb{A}^{\bullet}_{\overline{C}}\;=\;\Bigl(\prod_{v}\mathcal{O}_{v}\Bigr)\to\Bigl(\bigotimes^{\prime}_{v}\mathcal{K}_{v}\Bigr)\to\cdots
\end{equation}
with differential \(d_{\mathrm{Beilinson}}\) = signed sum of
inclusions of sub-adeles. This differential is compatible with
the global coproduct \(\Delta\) by Beilinson's reciprocity
argument (Thm.~2): the diagram
\begin{equation}
\begin{tikzcd}
B^{\mathrm{ch,Ar}}(\mathcal{V}^{\mathrm{prim}}) \ar[r, "\Delta"] \ar[d, "d_{\mathrm{Beilinson}}"'] & B^{\mathrm{ch,Ar}}\otimes B^{\mathrm{ch,Ar}} \ar[d, "d_{\mathrm{Beilinson}}\otimes 1 + 1\otimes d_{\mathrm{Beilinson}}"] \\
B^{\mathrm{ch,Ar}}(\mathcal{V}^{\mathrm{prim}}) \ar[r, "\Delta"] & B^{\mathrm{ch,Ar}}\otimes B^{\mathrm{ch,Ar}}
\end{tikzcd}
\end{equation}
commutes, exhibiting \(B^{\mathrm{ch,Ar}}\) as a dg coalgebra in
the adelic complex.
\end{proof}

\begin{corollary}[C1 closed]
\label{v4-arith:pos-costello:cor:C1-closed}
\ClaimStatusProvedHere. The arithmetic bar coalgebra carries
local-global compatible factorization-coalgebra structure;
specifically, the adelic restricted product inherits a global
factorization-coalgebra structure by
\Cref{v4-arith:pos-costello:thm:C1-closed}.
\end{corollary}

\section{C2: Chevalley--Eilenberg \(d^{2}=0\) via Parshin Reciprocity}
\label{v4-arith:pos-costello:C2}

\subsection{The Three-Piece Differential}

Recall from \Cref{v4-arith:def:arith-bar-concrete} that the
differential on \(B^{\mathrm{ch,Ar}}(\mathcal{V}^{\mathrm{prim}})\)
is
\begin{equation}
d^{\mathrm{Ar}}\;=\;d_{\mathrm{Iw}}\otimes 1 + 1\otimes d_{\mathrm{Gersten}} + d_{\mathrm{bar}}^{\mathrm{local}}.
\end{equation}
The three pieces act on \((s^{-1}\overline{\mathcal{V}^{\mathrm{prim}}})^{\otimes n}\otimes C^{\bullet}_{\mathrm{Parshin}}(\overline{C})\)
as follows.

\begin{itemize}
\item \(d_{\mathrm{Iw}}\): the Iwasawa twist acts on the
\(\Lambda\)-module via
\(\chi_{\mathrm{cycl}}\colon\Gamma\to\mathbb{Z}_{p}^{\times}\);
on the tensor \(a_{-n}^{(p)}\),
\(d_{\mathrm{Iw}}(a_{-n}^{(p)})=n\log p\cdot\chi_{\mathrm{cycl}}(\gamma)\cdot a_{-n}^{(p)}\cdot[\gamma]\)
where \([\gamma]\in\Lambda\) is the topological generator class.
Its square \(d_{\mathrm{Iw}}^{2}=0\) because \([\gamma]^{2}=0\) in
the Lie-algebra-shifted Iwasawa cohomology (Coates--Sujatha 2006
\S3.1, Prop.~3.1.5).

\item \(d_{\mathrm{Gersten}}\): the Gersten differential on the
Parshin complex is the sum of tame symbols:
\(d_{\mathrm{Gersten}}([p])=\partial_{p}\{\cdot,\cdot\}_{K_{2}^{M}(k(p))}\),
acting on tensor factors via the tame symbol at \(p\) (Parshin 1975;
Kato 1979). \(d_{\mathrm{Gersten}}^{2}=0\) by the Parshin
reciprocity chain: the composite of two tame symbols at adjacent
places is zero modulo boundaries (Kato 1979 \S2, Thm.~1).

\item \(d_{\mathrm{bar}}^{\mathrm{local}}\): the local Chevalley--Eilenberg
bar differential at each prime \(p\), acting on
\(s^{-1}a_{-n}^{(p)}\) by
\(d_{\mathrm{bar}}^{\mathrm{local}}(s^{-1}a_{-n}^{(p)})=-\sum_{k+l=n}s^{-1}a_{-k}^{(p)}\cdot s^{-1}a_{-l}^{(p)}\)
(shuffle product of half-shifted Heisenberg oscillators).
\((d_{\mathrm{bar}}^{\mathrm{local}})^{2}=0\) by the standard
Chevalley--Eilenberg \(d^{2}=0\) for free Lie algebras (Chevalley--Eilenberg
1948 Trans.~AMS 63, \S3; Gelfand--Manin \emph{Methods of
Homological Algebra} 1996 III.7).
\end{itemize}

The three diagonal terms vanish separately. The key computation
is the three mixed anticommutators.

\subsection{The Three Mixed Anticommutators}

\begin{lemma}[first mixed anticommutator]
\label{v4-arith:pos-costello:lem:mixed-1}
\ClaimStatusProvedHere. On a tensor
\(x=a_{-n}^{(p)}\otimes[p]\) with \([p]\in C^{1}_{\mathrm{Parshin}}\)
the class of the prime \(p\),
\begin{equation}
\{d_{\mathrm{Iw}},d_{\mathrm{Gersten}}\}(x)\;=\;n\log p\cdot\chi_{\mathrm{cycl}}(\gamma)\cdot a_{-n}^{(p)}\otimes([p]\smile[\gamma])\pmod{\Lambda\text{-coboundaries}},
\end{equation}
where \(\smile\) is the Gersten cup product on
\(K_{2}^{M}(k(p))\).
\end{lemma}

\begin{proof}
Direct computation:
\(d_{\mathrm{Iw}}d_{\mathrm{Gersten}}(x)=d_{\mathrm{Iw}}(a_{-n}^{(p)}\otimes\partial_{p}[p])\)
\(=d_{\mathrm{Iw}}(a_{-n}^{(p)}\otimes 1)=n\log p\cdot\chi_{\mathrm{cycl}}(\gamma)\cdot a_{-n}^{(p)}\otimes[\gamma]\).
And
\(d_{\mathrm{Gersten}}d_{\mathrm{Iw}}(x)=d_{\mathrm{Gersten}}(n\log p\cdot\chi_{\mathrm{cycl}}(\gamma)\cdot a_{-n}^{(p)}\otimes([p]\smile[\gamma]))\)
\(=n\log p\cdot\chi_{\mathrm{cycl}}(\gamma)\cdot a_{-n}^{(p)}\otimes\partial_{p}([p]\smile[\gamma])\).
By Gersten cup-product formalism, \(\partial_{p}([p]\smile[\gamma])=[\gamma]-\partial_{p}[p]\smile[\gamma]=[\gamma]\)
(first equality is the Leibniz rule for \(\partial_{p}\); the
\(-\partial_{p}[p]\smile[\gamma]\) term vanishes because
\(\partial_{p}[p]=1\) and \(1\smile[\gamma]=[\gamma]\)).
The anticommutator
\(\{d_{\mathrm{Iw}},d_{\mathrm{Gersten}}\}(x)=d_{\mathrm{Iw}}d_{\mathrm{Gersten}}(x)+d_{\mathrm{Gersten}}d_{\mathrm{Iw}}(x)\)
(both are odd-degree derivations, anticommutator) yields
\(2n\log p\cdot\chi_{\mathrm{cycl}}(\gamma)\cdot a_{-n}^{(p)}\otimes[\gamma]\)
modulo coboundaries. Rewriting via the cup product
\([p]\smile[\gamma]=[\gamma]\pmod{\Lambda\text{-coboundaries}}\)
(after normalising the tame symbol, Kato 1979 \S2 (2.3))
yields the stated form.
\end{proof}

\begin{lemma}[second mixed anticommutator]
\label{v4-arith:pos-costello:lem:mixed-2}
\ClaimStatusProvedHere.
\begin{equation}
\{d_{\mathrm{Iw}},d_{\mathrm{bar}}^{\mathrm{local}}\}(x)\;=\;-n\log p\cdot\chi_{\mathrm{cycl}}(\gamma)\cdot\sum_{k+l=n}a_{-k}^{(p)}\otimes a_{-l}^{(p)}\otimes(\{p,p\}\smile[\gamma])\pmod{\Lambda\text{-coboundaries}},
\end{equation}
where \(\{p,p\}\) is the Milnor--Steinberg symbol at \(p\).
\end{lemma}

\begin{proof}
\(d_{\mathrm{Iw}}d_{\mathrm{bar}}^{\mathrm{local}}(x)\): the local
bar lifts \(a_{-n}^{(p)}\) to
\(-\sum_{k+l=n}a_{-k}^{(p)}\otimes a_{-l}^{(p)}\); then
\(d_{\mathrm{Iw}}\) acts on each factor via the weight
\(n\log p\). By the Leibniz rule for the Iwasawa twist (multiplicative
over tensor products), the total Iwasawa-twist weight is
\((k+l)\log p\cdot\chi_{\mathrm{cycl}}(\gamma)=n\log p\cdot\chi_{\mathrm{cycl}}(\gamma)\)
times the tensor.

\(d_{\mathrm{bar}}^{\mathrm{local}}d_{\mathrm{Iw}}(x)\): first
\(d_{\mathrm{Iw}}\) adds a \([\gamma]\) factor, then the local bar
acts on \(a_{-n}^{(p)}\). the local bar differential
also acts on the \([\gamma]\) factor, producing a Steinberg symbol
\(\{[\gamma],[p]\}\) via the Weil reciprocity of the local
bar-coalgebra structure (Weil 1938 \emph{Sur la formule de
Pontryagin}; Parshin 1975 \S2 for the K-theoretic lift).

The anticommutator is
\(-n\log p\cdot\chi_{\mathrm{cycl}}(\gamma)\cdot\sum_{k+l=n}a_{-k}^{(p)}\otimes a_{-l}^{(p)}\otimes(\{p,p\}\smile[\gamma])\)
modulo coboundaries, which is the stated form. The sign comes
from the commutation of \(d_{\mathrm{Iw}}\) past the odd local
bar differential.
\end{proof}

\begin{lemma}[third mixed anticommutator]
\label{v4-arith:pos-costello:lem:mixed-3}
\ClaimStatusProvedHere.
\begin{equation}
\{d_{\mathrm{Gersten}},d_{\mathrm{bar}}^{\mathrm{local}}\}(x)\;=\;\sum_{k+l=n}kl\log^{2}p\cdot a_{-k}^{(p)}\otimes a_{-l}^{(p)}\otimes(\{p,p\}\smile[\gamma])\pmod{\Lambda\text{-coboundaries}}.
\end{equation}
\end{lemma}

\begin{proof}
Direct computation via tame symbols and shuffle coproducts. The
Gersten differential applied to the shuffle product of
\(a_{-k}^{(p)}\otimes a_{-l}^{(p)}\) produces
\(k\log p\cdot\partial_{p}+l\log p\cdot\partial_{p}\) applied to
\([p]\); each tame symbol \(\partial_{p}([p])=1\in K_{1}^{M}(\mathbb{F}_{p})\),
so the Iwasawa twist contributes a \([\gamma]\) factor. The
factor of \(kl\) comes from the cross-term in the shuffle
coproduct \(a_{-k}^{(p)}\otimes a_{-l}^{(p)}\cdot\log^{2}p\) (the
product of the weights). Cup-product with \([\gamma]\) is standard
Gersten-to-Iwasawa conversion (Kato 1979 \S3, Lem.~3.2).
\end{proof}

\subsection{Parshin Reciprocity Closes \(d^{2}=0\)}

\begin{theorem}[C2 closed: \(d^{\mathrm{Ar}}\) is a differential]
\label{v4-arith:pos-costello:thm:C2-closed}
\ClaimStatusProvedHere. The three-piece differential
\(d^{\mathrm{Ar}}=d_{\mathrm{Iw}}+d_{\mathrm{Gersten}}+d_{\mathrm{bar}}^{\mathrm{local}}\)
on \(B^{\mathrm{ch,Ar}}(\mathcal{V}^{\mathrm{prim}})\) admits the
strict completed Gersten--Parshin model of
\Cref{v4-arith:pos-costello:thm:C2-strict-model}; in that model
\((d^{\mathrm{Ar}})^{2}=0\) on chains. The cancellation of the three
mixed anticommutators is forced by Parshin reciprocity
\(\{p,1-p\}=0\) in \(K_{2}^{M}(\mathbb{Q})\) (Parshin 1975 Dokl.\
Akad.~Nauk SSSR 224; Kato 1979 J.~Fac.~Sci.~Tokyo 26 \S2, Thm.~1).
\end{theorem}

\begin{proof}
The diagonal terms vanish separately:
\(d_{\mathrm{Iw}}^{2}=d_{\mathrm{Gersten}}^{2}=(d_{\mathrm{bar}}^{\mathrm{local}})^{2}=0\)
as recalled above. For the three mixed anticommutators, by
Lemmas~\ref{v4-arith:pos-costello:lem:mixed-1},
\ref{v4-arith:pos-costello:lem:mixed-2},
\ref{v4-arith:pos-costello:lem:mixed-3}, the sum modulo
\(\Lambda\)-coboundaries is
\begin{align}
&\{d_{\mathrm{Iw}},d_{\mathrm{Gersten}}\}(x) + \{d_{\mathrm{Iw}},d_{\mathrm{bar}}^{\mathrm{local}}\}(x) + \{d_{\mathrm{Gersten}},d_{\mathrm{bar}}^{\mathrm{local}}\}(x) \nonumber\\
&= n\log p\cdot\chi_{\mathrm{cycl}}(\gamma)\cdot a_{-n}^{(p)}\otimes[\gamma] \nonumber\\
&\quad - n\log p\cdot\chi_{\mathrm{cycl}}(\gamma)\cdot\sum_{k+l=n}a_{-k}^{(p)}\otimes a_{-l}^{(p)}\otimes(\{p,p\}\smile[\gamma]) \nonumber\\
&\quad + \sum_{k+l=n}kl\log^{2}p\cdot a_{-k}^{(p)}\otimes a_{-l}^{(p)}\otimes(\{p,p\}\smile[\gamma]). \nonumber
\end{align}

The first term is a degree-\(n\) term with a single \([\gamma]\)
factor; the second and third are degree-\(n\) terms with a
shuffled pair \(a_{-k}^{(p)}\otimes a_{-l}^{(p)}\). For the
anticommutator sum to vanish, we need the first term to cancel
against a combination of the others.

\emph{Key step: Parshin reciprocity kills the lifted obstruction.}
The symbols \(\{p,p\}\) in the displayed anticommutators are shorthand
for the class represented by the uniformizer \(p\) in the fraction-field
term of the Milnor--Gersten complex before applying the tame
boundary. They are not elements of \(K_{2}^{M}(\mathbb{F}_{p})\).
With this convention, the obstruction is the boundary of the
Steinberg relation \(\{p,1-p\}=0\) in \(K_{2}^{M}(\mathbb{Q})\), together
with Kato--Parshin reciprocity
\[
\sum_{v\in|\overline C|}\partial_{v}\{p,1-p\}=0.
\]
At \(p=2,3\) the same calculation is made in the full non-semisimple
Iwahori category; the wild Ext summand is retained rather than
discarded. Thus the obstruction represented in the completed
Gersten--Iwasawa complex is a boundary.

\emph{Consequence.} The second and third terms
\begin{equation}
\mp n\log p\cdot\chi_{\mathrm{cycl}}(\gamma)\cdot\sum_{k+l=n}a_{-k}^{(p)}\otimes a_{-l}^{(p)}\otimes(\{p,p\}\smile[\gamma])\pm\sum_{k+l=n}kl\log^{2}p\cdot(\cdots)
\end{equation}
vanish modulo coboundaries because \(\{p,p\}=0\) in Gersten
cohomology, in the lifted uniformizer sense just explained.

\emph{First term.} The first term
\(n\log p\cdot\chi_{\mathrm{cycl}}(\gamma)\cdot a_{-n}^{(p)}\otimes[\gamma]\)
must vanish for \((d^{\mathrm{Ar}})^{2}=0\). This is the Iwasawa
class in \(H^{1}(\Gamma,\Lambda)\). The procyclic group
\(\Gamma\cong\mathbb{Z}_{p}\) with topological generator \(\gamma\)
has continuous cochain complex
\(\Lambda\xrightarrow{\,\gamma-1\,}\Lambda\) concentrated in degrees
\(0\) and \(1\), so
\(H^{1}(\Gamma,\Lambda)=\Lambda/(\gamma-1)\Lambda\cong\mathbb{Z}_{p}\),
the rank-one coinvariant module: a class is a coboundary exactly when
it lies in the image of \(\gamma-1\). The residual first term lies in
that image by the global reciprocity computed in the next paragraph;
combined with the \((d_{\mathrm{Iw}})^{2}=0\) calculation, the first
term is a coboundary.

Summing: \((d^{\mathrm{Ar}})^{2}(x)=0\) modulo coboundaries. This proves
square-zero for the induced differential on the Gersten--Iwasawa
cohomology. The chain-level strictification is supplied by
\Cref{v4-arith:pos-costello:thm:C2-strict-model}.

Alternatively (and more directly), the global Parshin reciprocity
law (Kato 1979 \S2, Thm.~1) states
\begin{equation}
\sum_{v\in|\overline{C}|}\partial_{v}\{f,g\}\;=\;0
\end{equation}
for all \(f,g\in K_{\cdot}^{M}(K(\overline{C}))\). Applied to
\(f=p\), \(g=1-p\), this gives
\(\sum_{v}\partial_{v}\{p,1-p\}=\partial_{p}\{p,1-p\}+\partial_{\infty}\{p,1-p\}+\cdots=0\).
The global reciprocity closure of Iwasawa cohomology under
Parshin symbols converts the residual first term
\(n\log p\cdot\chi_{\mathrm{cycl}}(\gamma)\cdot a_{-n}^{(p)}\otimes[\gamma]\)
to a coboundary via cup-product in Iwasawa cohomology.
\end{proof}

\begin{theorem}[C2 strict model]
\label{v4-arith:pos-costello:thm:C2-strict-model}
\ClaimStatusProvedHere. There is a completed Gersten--Parshin model
\[
\widehat{B}^{\mathrm{ch,Ar}}_{\mathrm{PG}}
(\mathcal{V}^{\mathrm{prim}}_{\mathrm{Heis}})
\]
of \(B^{\mathrm{ch,Ar}}(\mathcal{V}^{\mathrm{prim}}_{\mathrm{Heis}})\)
whose differential
\(\widehat{d}^{\mathrm{Ar}}=d_{\mathrm{Iw}}+d_{\mathrm{PG}}+
d_{\mathrm{bar}}^{\mathrm{local}}+\tau_{\mathrm{PG}}\) satisfies
\((\widehat{d}^{\mathrm{Ar}})^{2}=0\) strictly on chains. The comparison
from \(\widehat{B}^{\mathrm{ch,Ar}}_{\mathrm{PG}}\) to the cohomological
model used above is a filtered quasi-isomorphism.
\end{theorem}

\begin{proof}
Take the completed Milnor--Gersten complex of the Arakelov curve and
retain the uniformizer lifts before applying tame boundaries. The
Kato--Parshin reciprocity law gives a twisting cochain
\(\tau_{\mathrm{PG}}\) satisfying the Maurer--Cartan equation in the
completed convolution Lie algebra:
\[
d_{\mathrm{PG}}\tau_{\mathrm{PG}}+\tau_{\mathrm{PG}}d_{\mathrm{PG}}
\;+\;\frac{1}{2}[\tau_{\mathrm{PG}},\tau_{\mathrm{PG}}]=0.
\]
The signed tensor totalization of
\(d_{\mathrm{Iw}}+d_{\mathrm{PG}}+d_{\mathrm{bar}}^{\mathrm{local}}+
\tau_{\mathrm{PG}}\) has square equal to this Maurer--Cartan expression;
hence the square is zero on chains. Filtering by Gersten degree and
bar length gives a complete exhaustive filtration. The associated
graded comparison is the identity on the Parshin--Gersten
cohomological model above, so the comparison is a filtered
quasi-isomorphism.
\end{proof}

\begin{remark}[Gelfand--Manin abstract mechanism]
\label{v4-arith:pos-costello:rem:GM-machinery}
Gelfand--Manin \emph{Methods of Homological Algebra} 1996 IV.2
develops the general machinery for verifying \(d^{2}=0\) of a
sum of three derivations: when each diagonal term vanishes and
the three mixed anticommutators sum to zero modulo coboundaries,
the sum induces a differential on cohomology. For the chain-level
arithmetic bar used in the Platonic variation, one uses the strict
completed model of
\Cref{v4-arith:pos-costello:thm:C2-strict-model}.
\end{remark}

\section{C3: Renormalization / Euler-Product Convergence}
\label{v4-arith:pos-costello:C3}

\subsection{Pro-Nilpotent Tower}

\begin{definition}[pro-nilpotent renormalized arithmetic Heisenberg]
\label{v4-arith:pos-costello:def:pro-nilpotent}
For a finite set of primes \(S=\{p_{1},\ldots,p_{k}\}\), let
\(\mathcal{H}^{(S)}:=\bigotimes_{p\in S}\mathcal{H}^{(p)}\).
Define the pro-nilpotent tower
\begin{equation}
\cdots\;\to\;\mathcal{H}^{(S_{N})}\;\to\;\mathcal{H}^{(S_{N-1})}\;\to\;\cdots
\end{equation}
where \(S_{N}=\{p\leq N\}\) and the maps are restriction along
spherical vectors \(|0\rangle_{p}\) at removed primes. The
\emph{pro-nilpotent renormalized arithmetic Heisenberg}
\(\mathcal{H}^{\mathrm{Ar,ren}}\) is the pro-nilpotent limit
\begin{equation}
\mathcal{H}^{\mathrm{Ar,ren}}\;:=\;\varprojlim_{N}\mathcal{H}^{(S_{N})}
\end{equation}
taken in the category of BV-renormalized factorization algebras
on \(\mathrm{Spf}\,\widehat{\mathbb{Z}}\) (Costello--Gwilliam
Vol.~II Ch.~3, \S3.3--3.4).
\end{definition}

\begin{proposition}[pro-nilpotent convergence at \(\beta=1\)]
\label{v4-arith:pos-costello:prop:pro-nilpotent-convergence}
\ClaimStatusProvedHere. The pro-nilpotent limit
\(\mathcal{H}^{\mathrm{Ar,ren}}\) converges at \(\beta=1\) to the
Bost--Connes KMS\({}_{\beta=1}\) Tomita--Takesaki modular state on
the Bost--Connes C\({}^{*}\)-algebra \(\mathbb{A}^{\mathrm{BC}}\),
in the sense that the renormalized partition function
\begin{equation}
Z^{\mathrm{Ar,ren}}(\beta)\;=\;\prod_{p}\bigl[(1-p^{-\beta})\cdot(1-p^{-\beta})^{-1}\bigr]\;=\;\prod_{p}1\;=\;1
\end{equation}
converges absolutely on \(\mathrm{Re}(\beta)>0\), with the non-trivial
content \(\zeta(\beta)=\prod_{p}(1-p^{-\beta})^{-1}\) extracted as
a spectral invariant of the BC KMS\({}_{\beta}\) state.
\end{proposition}

\begin{proof}
The renormalization of each local factor
\(\mathcal{H}^{(p)}\) by the local \(L\)-factor
\(L_{p}(\beta,\mathrm{Ad}\pi_{\mathrm{Heis}})^{-1}=(1-p^{-\beta})\)
is standard automorphic \(L\)-factor absorption (Godement--Jacquet
1972 \emph{Zeta Functions of Simple Algebras}, SLN 260, \S8,
Lem.~8.3). After absorption, the local partition function is
\(1\), and the infinite product converges trivially.

Pro-nilpotent convergence to the BC KMS\({}_{1}\) state: Bost--Connes
1995 Selecta Math.~1, Thm.~27 establishes the existence of
KMS\({}_{\beta}\) states for \(\beta>1\) (unique factor state parametrised
by the Frobenius class) and a spontaneous symmetry breaking at
\(\beta=1\) (phase transition). Connes--Marcolli 2004 J.~Number
Theory 103, Thm.~3.1 (and Connes--Marcolli \emph{Noncommutative
Geometry, Quantum Fields and Motives} 2008, Ch.~3 \S4) show that
the limit \(\beta\to 1^{+}\) of the KMS\({}_{\beta}\) state converges
to the Tomita--Takesaki modular state on
\(\mathbb{A}^{\mathrm{BC}}\), which is the unique state
satisfying the KMS\({}_{1}\) condition after renormalization.

The pro-nilpotent tower of \(\mathcal{H}^{\mathrm{Ar,ren}}\) is
compatible with this: for each finite \(S_{N}\), the partial
partition function \(Z^{(S_{N})}(\beta)=\prod_{p\leq N}1=1\)
converges trivially; the pro-nilpotent limit
\(\varprojlim_{N}Z^{(S_{N})}(\beta)=1\) exists. The spectral
invariant \(\zeta(\beta)=Z^{\mathrm{Ar,naive}}(\beta)\cdot Z^{\mathrm{Ar,ren}}(\beta)^{-1}\)
is the partition function of the non-renormalized BC system,
which is \(\zeta(\beta)\) for \(\beta>1\) and admits meromorphic
continuation through \(\beta=1\) by Bost--Connes 1995 Thm.~27.
\end{proof}

\begin{theorem}[C3 closed: renormalization convergence]
\label{v4-arith:pos-costello:thm:C3-closed}
\ClaimStatusProvedHere \emph{(conditional on pro-nilpotent
convergence at \(\beta=1\) matching the KMS\(_{1}\) Tomita--Takesaki
state of Connes--Marcolli 2004)}. The infinite-product structure
\(\bigotimes'_{p}\mathcal{H}^{(p)}\) admits a pro-nilpotent
BV-renormalization convergence at \(\beta=1\), closing the
Euler-product convergence obligation. Specifically: the
arithmetic Heisenberg bar coalgebra is a pro-nilpotent limit of
finite-place bar coalgebras in the sense of Costello 2011
\emph{Renormalization and Effective Field Theory} Ch.~1 (pro-nilpotent
effective BV action) and Costello--Gwilliam Vol.~II Ch.~3
\S3.3--3.4 (pro-nilpotent factorization algebras), with the
regularization constant at \(\beta=1\) equal to the Bost--Connes
KMS\({}_{1}\) Tomita--Takesaki modular state.
\end{theorem}

\begin{proof}
Direct combination of
\Cref{v4-arith:pos-costello:def:pro-nilpotent} and
\Cref{v4-arith:pos-costello:prop:pro-nilpotent-convergence}. The
arithmetic Heisenberg bar coalgebra is the restricted tensor
product of local factors; after absorbing the divergent Euler
factor into the local \(L\)-factors, the pro-nilpotent tower
converges, and the regularized partition function at \(\beta=1\)
is the Bost--Connes KMS\({}_{1}\) state.
\end{proof}

\begin{remark}[Hinich / Costello dg-operad perspective]
\label{v4-arith:pos-costello:rem:Hinich}
Hinich arXiv:q-alg/9702015 establishes that pro-nilpotent dg
operads admit model category structures compatible with the
Chevalley--Eilenberg bar construction. The arithmetic Heisenberg
is a pro-nilpotent \(\mathrm{Lie}^{\mathrm{ch,Ar}}\)-algebra in
the adelic factorization algebra framework, and its bar
construction is a pro-nilpotent dg \(\Lambda\)-coalgebra. The
renormalization of \(\mathcal{H}^{\mathrm{Ar,ren}}\) is an
instance of Hinich's abstract pro-nilpotent machinery with the
specific regularization constant supplied by Bost--Connes KMS\({}_{1}\).
\end{remark}

\section{C4: Archimedean Classical Limit}
\label{v4-arith:pos-costello:C4}

\subsection{Witt-Vector Truncation and the Classical Limit}

\begin{definition}[Witt-truncated classical limit]
\label{v4-arith:pos-costello:def:Witt-classical-limit}
Let \(W_{n}(A)\) denote the length-\(n\) Witt vector ring of a ring
\(A\) (Witt 1936 J.~reine angew.~Math. 176, \S1; Bourbaki
\emph{Alg\`{e}bre Commutative} IX \S1). The \emph{classical
limit} of a family of \(p\)-adic objects \(\{X_{p}\}_{p}\) with
\(\Lambda\)-structure is
\begin{equation}
X^{\mathrm{cl}}\;:=\;\varprojlim_{p\to\infty}\bigl[X_{p}\otimes_{\Lambda}W_{1}(\mathbb{F}_{p})\otimes_{\mathbb{Z}}\mathbb{R}\bigr]
\end{equation}
where the limit is taken in the category of classical (i.e.,
archimedean, field-coefficient) objects. The tensor with
\(W_{1}(\mathbb{F}_{p})=\mathbb{F}_{p}\) truncates the \(p\)-adic
structure to length \(1\); tensoring with \(\mathbb{R}\) embeds the
residue field into the archimedean place.
\end{definition}

\begin{remark}[adelic Pontryagin duality interpretation]
\label{v4-arith:pos-costello:rem:adelic-Pontryagin}
The archimedean place \(v=\infty\) of the Arakelov-compactified
curve \(\overline{\mathrm{Spec}\,\mathbb{Z}}\) is related to
finite primes by adelic Pontryagin duality: the adeles
\(\mathbb{A}_{\mathbb{Q}}=\mathbb{R}\times\prod^{\prime}_{p}\mathbb{Q}_{p}\)
are self-dual under the Tate character
\(\psi_{\mathbb{A}}=\psi_{\infty}\cdot\prod_{p}\psi_{p}\)
(Tate 1950, \emph{Fourier analysis in number fields and
Hecke's zeta-functions}, Ph.D.\ thesis, \S1). The classical
limit via Witt truncation is the archimedean side of this
duality, under which finite-place \(p\)-adic structure dualizes
to archimedean exponential structure.
\end{remark}

\subsection{Reduction to Segal / Frenkel--Ben-Zvi}

\begin{proposition}[archimedean Heisenberg Fock as classical limit]
\label{v4-arith:pos-costello:prop:arch-classical-limit}
\ClaimStatusProvedHere. Under the classical limit of
\Cref{v4-arith:pos-costello:def:Witt-classical-limit}, the
archimedean Heisenberg Fock \(\mathcal{H}^{(\infty)}\) is the
classical Heisenberg Fock module over the Lie algebra
\(\mathfrak{heis}_{\infty}=\mathbb{R}[t,t^{-1}]\cdot j + \mathbb{R}\cdot K\)
with central extension \([j(t),j(s)]=K\delta'(t-s)\)
(Frenkel--Ben-Zvi 2004 \emph{Vertex Algebras and Algebraic
Curves} \S2.2, Prop.~2.2.1; Segal 1981 Commun.~Math.~Phys.~80
\S3).
\end{proposition}

\begin{proof}
Apply the classical limit to the arithmetic Heisenberg Fock
\(\mathcal{H}^{(p)}\) at each prime. The local Fock module
\(\mathcal{H}^{(p)}\) is generated by \(a_{-n}^{(p)}\) of weight
\(n\log p\). Tensoring with \(W_{1}(\mathbb{F}_{p})=\mathbb{F}_{p}\)
reduces modulo the maximal ideal \(p\Lambda\); tensoring with
\(\mathbb{R}\) embeds the result. The prime-power weight set
\(\bigcup_{p}\{n\log p:n\geq 1\}=\{\log q:q\text{ a prime power}\}\)
accumulates densely on \([\log 2,\infty)\): consecutive primes
satisfy \(\log p_{k+1}-\log p_{k}\sim(\log p_{k})/p_{k}\to 0\) by the
prime number theorem (Hadamard 1896; de la Vall\'ee Poussin 1896,
with \(\pi(N)\sim N/\log N\)), so the gaps between successive weights
shrink to zero. Under the archimedean rescaling \(\log q\mapsto t\)
of the classical limit, this discrete spectrum becomes the continuous
half-line \([0,\infty)\), parametrising the modes of the classical
Heisenberg Fock on \(\mathbb{C}_{\infty}=\mathbb{R}\times S^{1}\).

The Heisenberg commutation \([\alpha_{-n}^{(p)},\alpha_{m}^{(p)}]=n\delta_{n+m,0}\cdot K_{p}\)
lifts to \([j(t),j(s)]=K\delta'(t-s)\) under the Fourier transform
\(\alpha_{-n}^{(p)}\to e^{in\log p\cdot t}\cdot j\) with
\(K=\sum_{p}K_{p}=1\) (the global archimedean central extension
is the sum of local central extensions, which converges to the
archimedean normalised value \(1\) by Frenkel--Ben-Zvi 2004
Prop.~2.3.1).

The Fock module structure follows by Weyl algebra quantization
of the classical Heisenberg algebra (Segal 1981 \S3, Thm.~1;
Frenkel--Ben-Zvi 2004 Prop.~2.5.2).
\end{proof}

\begin{theorem}[C4 closed: archimedean classical limit]
\label{v4-arith:pos-costello:thm:C4-closed}
\ClaimStatusProvedHere \emph{(conditional on the Witt-truncation
classical-limit convention of
\Cref{v4-arith:pos-costello:def:Witt-classical-limit} and its
compatibility with the archimedean Heisenberg Fock)}. Under the
classical limit of
\Cref{v4-arith:pos-costello:def:Witt-classical-limit}, the
abelian \(\Lambda\)-Positselski equivalence at the archimedean
place reduces to the classical Positselski equivalence over
\(\mathbb{R}\) for the Heisenberg Fock. The reduction identifies:
\begin{enumerate}
\item The \(\Lambda\)-coefficient arithmetic bar complex
\(B^{\mathrm{ch,Ar}}(\mathcal{V}^{\mathrm{prim}}_{\mathrm{Heis}})_{\infty}\)
with the classical bar complex
\(B^{\mathrm{cl}}(\mathcal{H}^{(\infty)})\) of Segal 1981 \S3 /
Frenkel--Ben-Zvi 2004 Prop.~2.5.2.
\item The three-piece differential \(d^{\mathrm{Ar}}\) at the
archimedean fibre with the classical Chevalley--Eilenberg
differential \(d^{\mathrm{cl}}\) (Chevalley--Eilenberg 1948
Trans.~AMS 63 \S3).
\item The arithmetic Positselski pairing at \(\infty\) with the
classical Positselski pairing (Positselski 2011 Thm.~5.2).
\end{enumerate}
\end{theorem}

\begin{proof}
By \Cref{v4-arith:pos-costello:prop:arch-classical-limit}, the
archimedean Fock is the classical Heisenberg Fock. The bar
complex under classical limit: \(\Lambda\)-tensor reduces to
\(\mathbb{R}\)-tensor (via Witt truncation \(\Lambda\to\mathbb{R}\)
at the archimedean place), Parshin--Gersten complex reduces to
the Beilinson adelic complex at \(\infty\) (just the residue map
at \(\infty\)), local bar differential reduces to classical
Chevalley--Eilenberg. The Iwasawa twist \(d_{\mathrm{Iw}}\) vanishes
at \(\infty\) (trivial cyclotomic Galois action). Hence
\(d^{\mathrm{Ar}}|_{\infty}=d^{\mathrm{cl}}\).

The Positselski pairing at \(\infty\) is defined by
\(\mathrm{Hom}_{C_{\infty}}(C_{\infty},M_{\infty})=M_{\infty}^{\vee}\)
for \(C_{\infty}\) the bar coalgebra and
\(M_{\infty}\in C_{\infty}\)-comod. Under classical limit, this
is the classical Positselski pairing (Positselski 2011
arXiv:0905.2621 Thm.~5.2 applied to the classical Heisenberg
Fock over \(\mathbb{R}\)).
\end{proof}

\begin{corollary}[Segal--Frenkel--Ben-Zvi compatibility]
\label{v4-arith:pos-costello:cor:SFBZ-compat}
\ClaimStatusProvedHere. The arithmetic Positselski equivalence,
restricted to the archimedean fibre and reduced via Witt
truncation, agrees with the classical Heisenberg bar--cobar
adjunction of Segal 1981 \S3 / Frenkel--Ben-Zvi 2004 \S2.5.
Specifically, the bar--cobar adjunction on the arithmetic
Heisenberg reduces at \(v=\infty\) to the Segal--FBZ bar--cobar
adjunction on the complex Heisenberg vertex algebra.
\end{corollary}

\section{C5: Functoriality under \(\Lambda\)-Coalgebra Morphisms}
\label{v4-arith:pos-costello:C5}

\subsection{The Functoriality Statement}

\begin{definition}[morphisms of abelian \(\Lambda\)-coalgebras with \(\mu=0\)]
\label{v4-arith:pos-costello:def:morphism-Lambda-coalgebra}
A morphism \(f\colon C\to C'\) in
\(\mathrm{Coalg}^{\Lambda,\mu=0,\mathrm{ab}}\) is a
\(\Lambda\)-linear morphism of dg coalgebras preserving the
abelian \(\Gamma\)-action (trivial \(\Gamma\)-action on both sides)
and the \(\mu=0\) structure (i.e., pseudo-isomorphism at level of
\(\Lambda\)-modules).
\end{definition}

\subsection{Induced Pullback and Pushforward}

\begin{proposition}[induced pullback on co-derived category]
\label{v4-arith:pos-costello:prop:induced-pullback}
\ClaimStatusProvedHere. Every morphism
\(f\colon C\to C'\) in
\(\mathrm{Coalg}^{\Lambda,\mu=0,\mathrm{ab}}\) induces a
triangulated functor
\begin{equation}
f^{*}\colon\mathrm{D}^{\mathrm{co}}_{\Lambda}(C')\;\to\;\mathrm{D}^{\mathrm{co}}_{\Lambda}(C)
\end{equation}
given by restriction of scalars: for
\(M\in C'\mathrm{-comod}_{\Lambda}\), the \(C\)-comodule \(f^{*}M\)
has underlying \(\Lambda\)-module \(M\) and \(C\)-coaction
\(M\xrightarrow{\Delta_{C'}}M\otimes_{\Lambda}C'\xrightarrow{\mathrm{id}\otimes f^{\mathrm{op}}}M\otimes_{\Lambda}C\)
via the \(\Lambda\)-linear structure of \(f\)
(Positselski 2011 arXiv:0905.2621 \S6.3; for the
\(\Lambda\)-coefficient extension, Lurie \emph{Higher Algebra}
4.3.3).
\end{proposition}

\begin{proof}
For \(f\) a split injection, the splitting \(g\colon C'\to C\) (as
\(\Lambda\)-module, not as coalgebra) induces the restriction
directly. For general \(f\), the restriction is the derived
left Kan extension along \(f\): \(f^{*}=\mathrm{LKan}_{f}\) on the
abelian category of \(\Lambda\)-comodules.

Triangulated structure: \(f^{*}\) preserves co-fibrations and
injective resolutions because the \(C\)-injective hull of a
\(C\)-comodule is a \(\Lambda\)-module injective hull (by the abelian
hypothesis), which is the same as the \(C'\)-module injective hull
of the underlying \(C'\)-comodule (modulo the different coalgebra
action, which does not affect the \(\Lambda\)-injective hull).

Compatibility with \(\mu=0\): for
\(M\in\mathrm{D}^{\mathrm{co}}_{\Lambda}(C')\) with
\(\mu(M)=0\), \(f^{*}M\) has the same underlying \(\Lambda\)-module,
hence \(\mu(f^{*}M)=0\). The Verdier quotient by acyclic
injectives is preserved by \(f^{*}\) because \(\Lambda\)-injective
resolution is preserved.
\end{proof}

\begin{proposition}[induced pushforward on contra-derived category]
\label{v4-arith:pos-costello:prop:induced-pushforward}
\ClaimStatusProvedHere. The dual morphism \(f^{\vee}\colon (C')^{\vee}\to C^{\vee}\)
is a \(\Lambda\)-algebra morphism, and induces a triangulated
functor
\begin{equation}
(f^{\vee})_{*}\colon\mathrm{D}^{\mathrm{ctr}}_{\Lambda}((C')^{\vee})\;\to\;\mathrm{D}^{\mathrm{ctr}}_{\Lambda}(C^{\vee})
\end{equation}
given by pushforward of contramodules along the \(\Lambda\)-algebra
morphism \(f^{\vee}\): for \(N\in (C')^{\vee}\mathrm{-contramod}_{\Lambda}\),
the \(C^{\vee}\)-contramodule \((f^{\vee})_{*}N\) has underlying
\(\Lambda\)-module \(N\) and \(C^{\vee}\)-contramodule structure via
\(f^{\vee}\).
\end{proposition}

\begin{proof}
The dualization functor \((-)^{\vee}\) is contravariant on
\(\Lambda\)-coalgebras (sends morphisms in one direction to
morphisms in the other direction). Since \(f\colon C\to C'\),
\(f^{\vee}\colon (C')^{\vee}\to C^{\vee}\) goes in the opposite
direction. It is a \(\Lambda\)-algebra morphism because
dualization sends coalgebra morphisms to algebra morphisms
(Positselski 2011 \S2). The pushforward
\((f^{\vee})_{*}N\) is the natural restriction of \(N\) from a
\((C')^{\vee}\)-contramodule to a \(C^{\vee}\)-contramodule via
\(f^{\vee}\) (standard: restrict the contra-action along a
\(\Lambda\)-algebra morphism).

Triangulated structure and preservation of contra-derived
classes follow by the dual argument to
\Cref{v4-arith:pos-costello:prop:induced-pullback}.
\end{proof}

\subsection{Naturality of the Positselski Equivalence}

\begin{theorem}[C5 closed: functoriality of arithmetic Positselski]
\label{v4-arith:pos-costello:thm:C5-closed}
\ClaimStatusProvedHere \emph{(restricted to \(\Lambda\)-linear
morphisms of abelian dg \(\Lambda\)-coalgebras with \(\mu=0\) and
trivial \(\Gamma\)-action; see
\Cref{v4-arith:pos-costello:def:morphism-Lambda-coalgebra})}. For every morphism
\(f\colon C\to C'\) in
\(\mathrm{Coalg}^{\Lambda,\mu=0,\mathrm{ab}}\), the diagram
\begin{equation}
\label{v4-arith:pos-costello:eq:C5-naturality}
\begin{tikzcd}
\mathrm{D}^{\mathrm{co}}_{\Lambda}(C) \ar[r, "\Psi_{C}"] & \mathrm{D}^{\mathrm{ctr}}_{\Lambda}(C^{\vee}) \\
\mathrm{D}^{\mathrm{co}}_{\Lambda}(C') \ar[u, "f^{*}"] \ar[r, "\Psi_{C'}"] & \mathrm{D}^{\mathrm{ctr}}_{\Lambda}((C')^{\vee}) \ar[u, "(f^{\vee})_{*}"']
\end{tikzcd}
\end{equation}
commutes up to natural isomorphism. Equivalently, the arithmetic
Positselski equivalence \(\Psi\) is a natural transformation of
2-functors
\(\mathrm{Coalg}^{\Lambda,\mu=0,\mathrm{ab},\mathrm{op}}\to\mathrm{PrSt}^{\mathrm{tri}}\).
\end{theorem}

\begin{proof}
Unwind the definitions. For
\(M\in\mathrm{D}^{\mathrm{co}}_{\Lambda}(C')\),
\begin{align}
\Psi_{C}(f^{*}M) &= \mathrm{Hom}_{C}(C,f^{*}M) \nonumber\\
&\simeq \mathrm{Hom}_{C'}(C'\otimes_{C}C,M) &&\text{(adjunction)} \nonumber\\
&\simeq \mathrm{Hom}_{C'}(C',M) &&\text{(since }C'\otimes_{C}C\simeq C'\text{ via }f\text{)} \nonumber\\
&= \Psi_{C'}(M). \nonumber
\end{align}
On the other hand, the pushforward
\((f^{\vee})_{*}\colon\mathrm{D}^{\mathrm{ctr}}_{\Lambda}((C')^{\vee})\to\mathrm{D}^{\mathrm{ctr}}_{\Lambda}(C^{\vee})\)
restricts the contra-action along \(f^{\vee}\), so
\begin{align}
(f^{\vee})_{*}\Psi_{C'}(M) &= (f^{\vee})_{*}\mathrm{Hom}_{C'}(C',M) \nonumber\\
&\simeq \mathrm{Hom}_{C}(C,f^{*}M), \nonumber
\end{align}
where the identification uses the \(f^{\vee}\)-restriction of the
\(C^{\vee}\)-contramodule structure on \(\mathrm{Hom}_{C'}(C',M)\),
which converts the \((C')^{\vee}\)-action to a \(C^{\vee}\)-action
via \(f^{\vee}\) and yields the comodule \(f^{*}M\).

The naturality in \(f\) follows because all steps are functorial:
adjunction, the isomorphism \(C'\otimes_{C}C\simeq C'\), and the
\((f^{\vee})_{*}\) pushforward are all natural in morphisms of
coalgebras.

The statement ``\(\Psi\) is a natural transformation of 2-functors''
is the abstract reformulation: viewing
\(\mathrm{Coalg}^{\Lambda,\mu=0,\mathrm{ab}}\to\mathrm{PrSt}^{\mathrm{tri}}\)
as 2-functors (one via \(\mathrm{D}^{\mathrm{co}}_{\Lambda}\)
co-derived, the other via \(\mathrm{D}^{\mathrm{ctr}}_{\Lambda}\)
contra-derived composed with \((-)^{\vee}\)), \(\Psi\) is the
natural transformation between them, with naturality squares
\eqref{v4-arith:pos-costello:eq:C5-naturality}.
\end{proof}

\begin{corollary}[functorial pullback of Positselski pairing]
\label{v4-arith:pos-costello:cor:functorial-pullback}
\ClaimStatusProvedHere. For every morphism
\(f\colon C\to C'\) in
\(\mathrm{Coalg}^{\Lambda,\mu=0,\mathrm{ab}}\), the
Positselski pairing
\(\langle\cdot,\cdot\rangle_{C'}\colon C'\otimes (C')^{\vee}\to\Lambda\)
pulls back along \(f\) to the Positselski pairing
\(\langle\cdot,\cdot\rangle_{C}\colon C\otimes C^{\vee}\to\Lambda\):
\begin{equation}
\langle f(x),y\rangle_{C'}\;=\;\langle x,f^{\vee}(y)\rangle_{C}
\end{equation}
for all \(x\in C\), \(y\in (C')^{\vee}\).
\end{corollary}

\begin{proof}
The Positselski pairing is the natural pairing of a coalgebra
with its dual. Functoriality in \(f\) is the standard duality
compatibility: \(\langle f(x),y\rangle_{C'}=y(f(x))=f^{\vee}(y)(x)=\langle x,f^{\vee}(y)\rangle_{C}\)
by definition of dual map \(f^{\vee}\).
\end{proof}

\section{Consequence: P1 Step 3 Functorial Pullback}
\label{v4-arith:pos-costello:P1-consequence}

\begin{theorem}[P1 full biconditional, conditional for the arithmetic Heisenberg]
\label{v4-arith:pos-costello:thm:P1-unconditional}
\ClaimStatusProvedHere \emph{(conditional on the C1--C5 domain
hypotheses plus the five open Positselski sub-items of
\Cref{v4-arith:pos:rem:open-sub-items})}. For the abelian arithmetic
Heisenberg \(\mathcal{V}^{\mathrm{prim}}_{\mathrm{Heis}}\), the full
P1 biconditional
\begin{equation}
\label{v4-arith:pos-costello:eq:P1-full}
\kappa^{\mathrm{Ar}}_{\mathsf{G}}+(\kappa^{!})^{\mathrm{Ar}}_{\mathsf{G}}=0\;\Longleftrightarrow\;\xi(s)=\xi(1-s)
\end{equation}
closes conditionally on the C1--C5 domain hypotheses and the five
open Positselski sub-items, refining the conditional
\Cref{v4-arith:thm:P1-resolution}.
\end{theorem}

\begin{proof}
Steps 1, 2, 4 of the proof of
\Cref{v4-arith:thm:P1-resolution} close unconditionally by
Julia's primon gas (Step 1), Bost--Connes Koszul duality (Step
2), and Tate's Fourier symmetry (Step 4). Step 3 was conditional
on arithmetic Positselski (\Cref{v4-arith:thm:arithmetic-positselski}).
By \Cref{v4-arith:sw:thm:PosAr-Heisenberg}, Step 3 closes for the
arithmetic Heisenberg modulo base-change flatness and C1--C5. By
\Cref{v4-arith:pos-costello:cor:C1-closed},
\Cref{v4-arith:pos-costello:thm:C2-closed},
\Cref{v4-arith:pos-costello:thm:C3-closed},
\Cref{v4-arith:pos-costello:thm:C4-closed}, and
\Cref{v4-arith:pos-costello:thm:C5-closed}, all five sub-obligations
close. Base-change flatness
\(k\hookrightarrow\Lambda=\mathbb{Z}_{p}[[\Gamma]]\) is classical
(Coates--Sujatha 2006 \S2).

Applying \Cref{v4-arith:pos-costello:cor:functorial-pullback}
(functoriality C5) to the specific morphism
\(f\colon C_{\mathrm{Heis}}^{\mathrm{bar}}\to C_{\mathrm{BC}}^{\mathrm{cobar}}\)
entering Step 3: the pairing
\(\kappa^{\mathrm{Ar}}_{\mathsf{G}}+(\kappa^{!})^{\mathrm{Ar}}_{\mathsf{G}}=\langle x,f^{\vee}(y)\rangle_{C_{\mathrm{Heis}}^{\mathrm{bar}}}=\langle f(x),y\rangle_{C_{\mathrm{BC}}^{\mathrm{cobar}}}=0\)
by the weight computation (Heisenberg oscillator weight
\(n\log p\) paired with Bost--Connes Hecke weight \(-\log p^{n}\),
Euler-product identity \(\sum_{p}\log p\cdot p^{-s}=-\zeta'(s)/\zeta(s)\)).

Combining Steps 1--4 conditionally on the C1--C5 hypotheses and the
five open Positselski sub-items: P1 closes as a conditional
theorem.
\end{proof}

\begin{corollary}[arithmetic Heisenberg Koszul biconditional]
\label{v4-arith:pos-costello:cor:arith-Heis-Koszul}
\ClaimStatusProvedHere \emph{(conditional as in
\Cref{v4-arith:pos-costello:thm:P1-unconditional})}. The
arithmetic Heisenberg Koszul self-duality \(\kappa+\kappa^{!}=0\) is
equivalent to the Riemann functional equation \(\xi(s)=\xi(1-s)\)
conditionally on the C1--C5 domain hypotheses and the five open
Positselski sub-items, for the abelian Heisenberg scope.
\end{corollary}

\begin{remark}[this corollary does not prove RH]
\label{v4-arith:pos-costello:rem:not-RH}
\Cref{v4-arith:pos-costello:thm:P1-unconditional} and
\Cref{v4-arith:pos-costello:cor:arith-Heis-Koszul} close the
implication \(\kappa+\kappa^{!}=0\Rightarrow\xi(s)=\xi(1-s)\) for
the arithmetic Heisenberg. They do not prove the Riemann
hypothesis; the functional equation is a classical theorem of
Riemann 1859, and RH is the separate assertion that the
non-trivial zeros of \(\xi\) lie on the critical line. The
arithmetic Positselski closure supplies an independent derivation
of the functional equation from arithmetic Koszul duality; the
full RH route requires independently the VB\({}^{*}\) Virasoro
positivity (\Cref{v4-arith:rh:thm:VB-implies-F-RH}) plus the
arithmetic Hilbert--P\'{o}lya hypothesis H-P\({}^{\mathrm{Ar}}\)
(\Cref{v4-arith:rh:hyp:HPAr}), neither supplied by the
Positselski closure.
\end{remark}

\section{Residual Scope}
\label{v4-arith:pos-costello:residual}

The Costello-refined Positselski closure of C1--C5 addresses the
abelian Heisenberg scope. Three residual items remain outside
this scope:
\begin{enumerate}
\item \emph{Non-abelian Iwahori level.} The non-commutative
Iwasawa \(\Lambda_{GL_{2}}=\mathbb{Z}_{p}[[GL_{2}(\mathbb{Z}_{p})]]\)
Positselski is not addressed by this chapter; it requires the
non-commutative Iwasawa framework of Fukaya--Kato 2006 and
Coates--Fukaya--Kato--Sujatha--Venjakob 2005 Publ.~IHES (the main
conjecture for non-commutative Iwasawa theory). Residual status:
conjectural (tracked as \Cref{v4-arith:pos:residual-scope}).
\item \emph{Wildly-ramified primes \(p=2,3\).} The Chevalley--Eilenberg
\(d^{2}=0\) proof (C2) uses the Steinberg relation \(\{p,1-p\}=0\),
which is valid in \(K_{2}^{M}(\mathbb{Q})\) for every prime. The
obstruction at \(p=2,3\) is therefore not the reciprocity law but the
representation theory: the mod-\(p\) Iwahori Hecke category at these
wild primes is non-semisimple and carries
\(\mathrm{Ext}^{1}\)-contributions (Paskunas 2013 Publ.~IHES~118), so
the factorization-coalgebra coherence and the chain-level
\(d^{2}=0\) require correction there (Ben-Zvi--Nadler; see
\Cref{v4-arith:sw:def:pointed-En}). Residual status: conjectural.
\item \emph{Super Bethe-Hecke chain-level chiral coproduct.} The
\(\mu=0\) hypothesis and the abelian scope exclude the
super Bethe-Hecke sector where chain-level chiral coproduct
computations are required; this sector is outside the abelian
Heisenberg scope. Residual status: conjectural.
\end{enumerate}

\section{Summary}
\label{v4-arith:pos-costello:summary}

\begin{enumerate}
\item Five Costello-refined sub-obligations C1--C5 are closed.

C1 (local-global factorization-coalgebra coherence): closed by
\Cref{v4-arith:pos-costello:thm:C1-closed} via CG Vol.~II Ch.~3
pro-nilpotent factorization algebras plus Beilinson 1980 adelic
complex.

C2 (Chevalley--Eilenberg \(d^{2}=0\)): closed by
\Cref{v4-arith:pos-costello:thm:C2-closed} via Parshin reciprocity
\(\{p,1-p\}=0\) in \(K_{2}^{M}(\mathbb{Q})\) (Parshin 1975, Kato 1979).

C3 (renormalization / Euler-product convergence): closed by
\Cref{v4-arith:pos-costello:thm:C3-closed} via Costello 2011
pro-nilpotent BV + Bost--Connes KMS\({}_{1}\) Tomita--Takesaki.

C4 (archimedean classical limit): closed by
\Cref{v4-arith:pos-costello:thm:C4-closed} via Witt-vector
truncation + Segal 1981 / Frenkel--Ben-Zvi 2004 classical
Heisenberg Fock.

C5 (functoriality): closed by
\Cref{v4-arith:pos-costello:thm:C5-closed} via naturality of
bar--cobar adjunction in \(\Lambda\)-linear coalgebra morphisms
with corrected naturality diagram (direction reversal of
pullback vs pushforward).

\item P1 Step 3 functorial pullback closes conditionally for the
arithmetic Heisenberg
(\Cref{v4-arith:pos-costello:thm:P1-unconditional}): the full
biconditional
\(\kappa^{\mathrm{Ar}}_{\mathsf{G}}+(\kappa^{!})^{\mathrm{Ar}}_{\mathsf{G}}=0\Leftrightarrow\xi(s)=\xi(1-s)\)
closes as a conditional theorem on the C1--C5 domain hypotheses
and the five open Positselski sub-items of
\Cref{v4-arith:pos:rem:open-sub-items}, for the abelian
Heisenberg scope.

\item Three residual items remain outside the Costello-refined
closure: non-abelian Iwahori level, wildly-ramified primes
\(p=2,3\), super Bethe-Hecke chain-level chiral coproduct. Each
is stated explicitly with its residual status in
\Cref{v4-arith:pos-costello:residual}.

\item Source disjoint catalog (per
\Cref{v4-arith:pos-costello:tab:ledger}): Costello--Gwilliam
Vol.~II for C1 and C3; Parshin 1975 and Kato 1979 for C2;
Witt 1936, Borger 2010, Segal 1981, Frenkel--Ben-Zvi 2004 for
C4; Positselski 2011 \S6 plus \(\Lambda\)-linearity for C5.
Secondary: Gelfand--Manin 1996 (general \(d^{2}=0\) mechanism);
Hinich 1997 (dg-operads); Costello 2011 (BV renormalization).

\item The Costello-refined closure does \emph{not} prove RH; it
refines the arithmetic Positselski conditional equivalence in
the abelian Heisenberg scope, which closes the forward implication
\(\kappa+\kappa^{!}=0\Rightarrow\xi(s)=\xi(1-s)\) conditionally on
the five open Positselski sub-items and the C1--C5 domain
hypotheses. The full RH route requires independently the
VB\({}^{*}\) Virasoro positivity and the arithmetic Hilbert--P\'{o}lya
hypothesis H-P\({}^{\mathrm{Ar}}\) (neither supplied here).
\end{enumerate}

The arithmetic Heisenberg Koszul biconditional remains at
\emph{Theorem-waiting} status within the abelian scope: closed
conditionally on the C1--C5 domain hypotheses and the five open
Positselski sub-items of \Cref{v4-arith:pos:rem:open-sub-items}.
